\metatitle{Real-rooted polynomials from set partitions}

\metadescription{Real-rootedness and interlacing for set partitions,
Stirling permutations, Jacobi--Stirling permutations, and related words.}
\metakeywords{real-rooted polynomial, set partition, Stirling permutation,
Jacobi-Stirling permutation, quasi-Stirling permutation, Sturm sequence}


\section[sturmSetPartitions]{Set partitions}



\begin{example*}[Touchard polynomials (blocks in set partitions), with proof]\label{touchardPolynomial}

The \defin{Touchard polynomials}\label{touchardPolynomialDefinition}
may be defined via
\[
T_n(x) \coloneqq \sum_{P \in SP(n)} x^{\mathrm{blocks}(P)}
= \sum_{k=1}^n S(n,k) x^k,
\]
where $SP(n)$ is the set of set partitions of $[n]$, and we count the number
of blocks.
Note that the coefficients are Stirling numbers of the second type.

Wilf gives the recurrence \cite[pp.~20, 139]{Wilf1994}
\[
  T_n(x) = x \cdot  T_{n-1}(x) + x\cdot T'_{n-1}(x).
\]
We know $T_{n-1} \interl x T_{n-1}$. Moreover, $T_{n-1} \interl x T'_{n-1}$,
so by \cite{Wagner1992}, we have $T_{n-1} \interl x T_{n-1} +  x T'_{n-1}$.
But this is exactly $T_{n-1} \interl T_n$.

Another way to prove real-rootedness is by interpreting \(T_n(x)\) as a
\hyperref[touchardAsMatching]{matching polynomial}.
\end{example*}


\begin{example*}[Set partitions without singleton blocks]

Define \(D_n(t)=\sum_{\pi \in SP_2(n)} t^{\mathrm{blocks}(\pi)}\),
where $SP_2(n)$ is the set of set partitions of $[n]$ whose blocks all have
size at least $2$.
Then
\[
 D_n(t) = t\left( D'_{n-1}(t) + (n-1) D_{n-2}(t) \right)
\]
and $D_{n-1}(t)$ and $D_n(t)$ have interlacing roots, see
\cite[Thm.~1]{BonaHezo2016}.
Note that the degree of $D_n(t)$ is $\lfloor n/2 \rfloor$.


The proof is short by using Wagner's lemma and induction.
The base case is immediate; assume $D_{n-2} \interl D_{n-1}$ for some
$n \geq 2$.
We then know that
\[
D'_{n-1} \interl D_{n-1} \text{ and } (n-1)D_{n-2} \interl D_{n-1}.
\]
Wagner's lemma now gives that
\[
D'_{n-1} + (n-1)D_{n-2} \interl D_{n-1}
\implies
D_{n-1} \interl t \left( D'_{n-1} + (n-1)D_{n-2} \right),
\]
and this is precisely $D_{n-1} \interl D_n$, which completes the induction step.

This result also follows from the general machinery in
\cite[Subsec.~3.4]{Hitczenko2024x}.
A recursion for set partitions into blocks of size at least $s$ is also mentioned
there (originally due to \name{Louis Comtet})
but it does not immediately imply real-rootedness.

\end{example*}



\begin{example*}[Type $B$ set partitions]

\name{David Wang} considers a type $B$ version, and more generally a colored
version, of counting blocks in set partitions \cite{Wang2014}.
The polynomials considered satisfy the (general) recursion
\[
  T_n(x) = (x+c) \cdot  T_{n-1}(x) + mx\cdot T'_{n-1}(x),
\]
where $m$ and $c$ are positive integers.

\name{Per Alexandersson}, \name{Fufa Beyene}, and \name{Roberto Mantaci}
study type $B$ set partitions without a zero block
\cite{AlexanderssonBeyeneMantaci2026x}.  If $b_{n,k}$ counts the
merging-free ones with $k$ blocks and
$M_n(t)=\sum_k b_{n,k}t^k$, then
\[
 M_n(t)=2tM'_{n-1}(t)+2(n-2)tM_{n-2}(t),\qquad n\geq2,
\]
with $M_0(t)=1$ and $M_1(t)=t$.
The polynomials $M_n(t)$ are real-rooted and consecutive members interlace
\cite[Prop.~27 and Cor.~28]{AlexanderssonBeyeneMantaci2026x};
the paper also proves real-rootedness for the block polynomials of type $B$
partitions that are both merging-free and separated
\cite[Thm.~33]{AlexanderssonBeyeneMantaci2026x}.
Its flattened signed-permutation descent polynomial is gamma-positive and
real-rooted \cite[Thm.~59 and Cor.~60]{AlexanderssonBeyeneMantaci2026x}.

\end{example*}


\begin{example*}[Ordered partitions]
Simion considers any multiset $M$ and the number $O(M,k)$ of compositions into
exactly $k$ parts \cite{Simion1984}.  The polynomial
$\sum_k O(M,k)x^k$ is real-rooted.

In particular, this implies the real-rootedness of the Touchard polynomials.
\end{example*}


\subsection[sturmOther]{Stirling and quasi-Stirling permutations}


\begin{example*}[Stirling permutations]

The descent number of a \hyperref[permutationsStirling]{Stirling permutation}
\(\pi \in Q_n\) is the cardinality of
\[
 \{ i : \pi_i \gt \pi_{i+1} \} \cup \{2n\}.
\]
The \defin{Stirling permutation descent polynomials}
\label{stirlingPermutationDescentPolynomial} satisfy the recurrence
\begin{equation}
 P_n(t) = (2n+1)t \cdot P_{n-1}(t) + t(1-t) P'_{n-1}(t),
\end{equation}
see \cite[Eq.~13]{Carlitz1965} and \cite[Thm.~2.1]{GesselStanley1978}.
The coefficients can be found in \oeis{A008517}; the first few values are
presented below.
\begin{array}{rrrrrr}
\toprule
 1  \\
 1 & 2  \\
 1 & 8 & 6 \\
 1 & 22 & 58 & 24 \\
 1 & 52 & 328 & 444 & 120 \\
 1 & 114 & 1452 & 4400 & 3708 & 720 \\
\bottomrule
\end{array}
These numbers appear in the Hilbert series of the Whitehouse complex;
see \cite{Readdy2005}.  These polynomials are the
\hyperref[hStarPolynomial]{$h^*$-polynomials} of the
\hyperref[orderPolytope]{order polytope} defined by the inequalities
\[
 y_1 \leq y_2 \leq \dotsb \leq y_n \qquad 0 \leq y_j \leq x_j \leq 1 \text{ for } j=1,\dotsc,n.
\]
\name{Shi-Mei Ma}, \name{Jun Ma}, \name{Jean Yeh}, and
\name{Yeong-Nan Yeh} prove $e$-positivity for multivariate $k$th-order
Eulerian polynomials attached to $k$-Stirling permutations
\cite{MaMaYehYeh2021x}.  They also give tree interpretations for the
coefficients in related gamma-type expansions.
For $n=20$, the linear term of the Ehrhart polynomial (of degree $40$) has
negative coefficient $-22354369153/6983776800$.

\name{Miklós Bóna} proved that these have only real zeros \cite{Bona2009}.
The Ma--Wang theorem also gives real-rootedness and interlacing.

\name{Jim Haglund} and \name{Mirko Visontai} \cite[Thm.~3.3]{HaglundVisontai2012}
show that a multivariate generalization is stable.
\end{example*}

\begin{theorem*}[Dey's interlacing criterion, \cite{Dey2023InterlacingZeroes}]
\label{deyInterlacingCriterion}
Let $(F_n(x))_{n\geq 0}$ be a sequence of polynomials satisfying
\[
  F_n(x)=\bigl(\alpha_1(n)x+\alpha_2\bigr)F_{n-1}(x)
  +\alpha_3x(1-x)F'_{n-1}(x),
\]
where $\alpha_2,\alpha_3\in\setN\cup\{0\}$. Suppose that each
$F_n(x)$ is real-rooted with distinct nonpositive zeros, and that
some integer $k$ satisfies, for all $n\geq k$,
\[
  2\alpha_1(n)-n\alpha_3>0,
  \qquad
  n\alpha_1(n)\alpha_3+3n\alpha_2\alpha_3
  -4\alpha_1(n)\alpha_2>0.
\]
Then $F_{n-1}(x)\interl F_n(x)$ for all $n\geq k$.
\end{theorem*}

Dey applies this criterion to the type $A$ and type $B$ Eulerian
polynomials, to the Stirling permutation polynomials above, and to
\hyperref[restrictedEulerianPolynomial]{restricted Eulerian polynomials}.
Thus it is a useful way to upgrade an
already known real-rootedness proof to a consecutive interlacing
statement when the recurrence has the form above.


\subsubsection[jacobiStirlingDescentRealRooted]{Jacobi--Stirling descent polynomials}

For $k\geq1$, let
\[
 M_k=\{\overline{1},1,1,\overline{2},2,2,\dotsc,\overline{k},k,k\},
 \qquad
 \overline{1}<1<\overline{2}<2<\dotsb<\overline{k}<k.
\]
For $S\subseteq[k]$, delete the barred letters $\overline{j}$ with $j\in S$.
A \defin{Jacobi--Stirling permutation} of the resulting multiset is a word
in which every entry between the two unbarred copies of $j$ is larger than
$j$.  Let $\operatorname{JSP}_{k,S}$ be the set of these words, and define
\[
 A_{k,S}(t)\coloneqq\sum_{w\in\operatorname{JSP}_{k,S}}t^{\des(w)},
 \qquad
 A_{k,i}(t)\coloneqq
 \sum_{\substack{S\subseteq[k]\\|S|=i}}A_{k,S}(t).
\]
Here $\des$ counts internal descents.  In particular, $A_{k,k}(t)$ is the
second-order Eulerian polynomial with the final descent omitted, and
$[t^j]A_{k,k}(t)$ is entry $(k,j+1)$ of \oeis{A008517}.

\name{Ira Gessel}, \name{Zhicong Lin}, and \name{Jiang Zeng} observed that
every $A_{k,S}(t)$ is real-rooted and conjectured the same for every
$A_{k,i}(t)$ \cite[Sec.~5 and Conj.~15]{GesselLinZeng2012}.  Their
normalization has an additional factor of $t$.

\begin{theorem*}[\name{Shi-Mei Ma} and \name{Ming-Xin Wang},
\cite[Thm.~2 and Cor.~2]{MaWang2026x}]
Let $k\geq1$ and
\[
 i\in\{1,2,3,k-2,k-1\}\cap[k].
\]
Every nonzero nonnegative linear combination of the polynomials
$A_{k,S}(t)$ with $|S|=i$ has only simple negative zeros.  In particular,
$A_{k,i}(t)$ has only simple negative zeros.
\end{theorem*}

The cases $i=0$ and $i=k$ were already known; the conjecture remains open
for the remaining values of $i$.

\begin{theorem*}[\cite[Thms.~3 and 5]{MaWang2026x}]
Write $A_{k,i}(t)=\sum_{j=0}^{d}a_jt^j$, where $k\geq2$ and
$d=2k-i-1$.  Then
\[
 a_j<a_{d-j}\quad(0\leq j<d/2),
 \qquad
 a_0<a_1<\dotsb<a_{\lfloor d/2\rfloor}.
\]
The same inequalities hold for every $A_{k,S}(t)$, with
$d=2k-|S|-1$.
\end{theorem*}

\subsubsection[quasiStirlingDescentRealRooted]{Quasi-Stirling descent polynomials}

\begin{example*}[Quasi-Stirling permutations]
A \defin{quasi-Stirling permutation}\label{quasiStirlingPermutation}
of size $n$ is a permutation of the multiset
$\{1,1,2,2,\dotsc,n,n\}$ avoiding the patterns $1212$ and $2121$.
Let $\widetilde{Q}_n$ be the set of quasi-Stirling permutations of size $n$,
and set
\[
 \widetilde{Q}_n(t) \coloneqq \sum_{\pi \in \widetilde{Q}_n} t^{\des(\pi)},
\]
where the final position is counted as a descent, as above.

\name{Sergi Elizalde} \cite[Thm.~3.2]{Elizalde2021} proved that
$\widetilde{Q}_n(t)$ has real, distinct, and nonpositive roots.
More generally, for $1 \leq r \leq n$, let
\[
 A_{n,r}(t) \coloneqq \sum_{\pi \in \symS_n} t^{\exc_r(\pi)}, \qquad
 J_{n,r}(t) \coloneqq \sum_{\pi \in \operatorname{Inj}([n-r],[n])} t^{\exc(\pi)},
\]
where $\exc_r(\pi) \coloneqq |\{i:\pi_i \geq i+r\}|$.
Then $A_{n,r}(t)$ and $J_{n,r}(t)$ also have real, distinct, and nonpositive
roots.
The quasi-Stirling case follows from the relation
\[
 \widetilde{Q}_n(t)=t J_{2n,n+1}(t).
\]
Elizalde also proves asymptotic normality for the descent distribution and
extends the analysis to a one-parameter family generalizing
\(k\)-Stirling permutations.
\end{example*}
